\section*{第 \ref{ch:b2:himalaya-andes} 章——碰撞造山带与安第斯型造山带的比较}

\begin{solution}{exo:b2:himalaya-andes:1}
(a) $r = h\rho_c/(\rho_m - \rho_c) = 3.0\times2700/600 = \qty{13.5}{km}$. (b) $T = 35 + 3.0 +
13.5 = \qty{51.5}{km}$，或 $35 + 5.5\times3.0$。(c) $h = \Delta T(\rho_m - \rho_c)/\rho_m =
10\times600/3300 = \qty{1.8}{km}$.
\end{solution}

\begin{solution}{exo:b2:himalaya-andes:2}
(a) 印度–恒河前陆坳陷；主前缘逆冲断层；次喜马拉雅（西瓦利克）；主边界
逆冲断层；小喜马拉雅；主中央逆冲断层；大喜马拉雅岩系；藏南
拆离系；特提斯喜马拉雅岩系；随后是印度河–雅鲁藏布缝合带和亚洲
地壳。(b) 藏南拆离系。(c) 变质级向北升高：未变质
\omterm{def:b2:alpine-orogen:foreland}{磨拉石}，中级变质的小喜马拉雅，高级变质且\omterm{def:b2:mantle-melting:melting}{部分熔融}的大喜马拉雅
岩系；在 STDS 之上，变质级骤降至未变质或弱变质的
特提斯沉积岩，因为正断层已移除其间的岩石。
\end{solution}

\begin{solution}{exo:b2:himalaya-andes:3}
(a) $196.5\ \text{m}/9200\ \text{年} = \qty{0.02136}{m}$ 每年，约 \qty{21.4}{mm} 每年。
(b) 比 GPS 的 \qty{17.8}{mm} 每年多 \qty{20}{\percent}，但量级相同：
一个速率是九千年的平均，另一个是二十年的平均。
(c) 整条山脉的缩短几乎全部集中在前缘：MBT、MCT 及
MFT 以北的其他断层在该经度上实际已不活动；它们都根植于
同一条 MHT，它出露于 MFT 处。
\end{solution}

\begin{solution}{exo:b2:himalaya-andes:4}
(a) 碰撞型：缝合带标志着已消失的海洋。(b) 安第斯型：火山弧需要
俯冲的板片。(c) 通常为碰撞型：加厚、受热地壳的熔融产物（它们也
出现于加厚的安第斯地壳中，但那里以弧\omterm{def:b2:magma-emplacement:magma}{岩浆}为主）。(d) 安第斯型。
(e) 两者皆有：西瓦利克带和次安第斯前陆。(f) 两者皆有：喜马拉雅的 MHT，
安第斯的俯冲界面。
\end{solution}

\begin{solution}{exo:b2:himalaya-andes:5}
(a) $17.8/45 = \qty{40}{\percent}$ 在东段，$20.5/45 = \qty{46}{\percent}$ 在西段：
约一半。(b) 一部分进入西藏及亚洲其余地区的增厚与断层，
另一部分进入沿大型走滑断层向东挤出的块体。(c) $45 - 17.8 =
\qty{27.2}{mm}$ 每年，或在 \qty{10}{Ma} 内 \qty{272}{km}；$45 - 20.5 = \qty{24.5}{mm}$ 每
年，或 \qty{245}{km}。
\end{solution}

\begin{solution}{exo:b2:himalaya-andes:6}
(a) $r = 4.5\times3.7 = \qty{16.65}{km}$; $\Delta F = 2700\times9.8\times3700\times(35\,000 +
1850 + 8325) = \qty{4.42e12}{N.m^{-1}}$。(b) 西藏：\qty{6.12e12}{N.m^{-1}}；阿尔蒂普拉诺
占其 \qty{72}{\percent}。(c) 取 $r = 4.5h$，$\Delta F = \rho_c g\,(Hh + 2.75h^2)$：
当 $h \ll H$ 时，$Hh$ 项占主导，$\Delta F$ 与 $h$ 成正比；而对于高原
，$h^2$ 项增长更快。高原每再升高一千米，边缘处所需的推力便越来越大，
直到可用的力再也无法抬升高原；此时高原
转为向外铺展，而不再隆起。
\end{solution}

\begin{solution}{exo:b2:himalaya-andes:7}
(a) $2700\times9.8\times8848.86/2 = \qty{1.17e8}{Pa} = \qty{117}{MPa}$。(b) 在
\qty{100}{MPa}：$0.85\times100 = \qty{85}{MPa}$；在 \qty{200}{MPa}：$0.85\times200 = 50 +
0.6\times200 = \qty{170}{MPa}$，两支结果吻合。(c) $h = 2\sigma_d/(\rho g) =
2\times1.70\times10^8/(2700\times9.8) = \qty{12.8}{km}$：数量级上
与最高峰相当，且并不高出太多。
\end{solution}

\begin{solution}{exo:b2:himalaya-andes:8}
(a) $11\times111.19\times\cos\ang{30} = \qty{1059}{km}$. (b) $9\times10^{-3}/1.059\times10^{6}
= \num{8.5e-9}$ 每年（7 和 \qty{11}{mm} 每年时，介于 \num{6.6e-9} 与 \num{10.4e-9} 之间），
即 $8.5\times10^{-9}/3.156\times10^{7} = \qty{2.7e-16}{s^{-1}}$。(c) 两段范围
重叠：$(12.4 \pm 4.7) \times 10^{-9}$ 覆盖 7.7 至 \num{17.1e-9}。位于
尼泊尔以北的南部高原伸展速率与整个山脉相当，甚至略快。
\end{solution}

\begin{solution}{exo:b2:himalaya-andes:9}
(a) $0.9/6.6 = \qty{14}{\percent}$. (b) $M_0 = 10^{1.5\times8.5 + 9.1} = 10^{21.85} =
\qty{7.1e21}{N.m}$。(c) $7.1\times10^{21}/6.6\times10^{19} = 107$：尼泊尔下方约
每 110 年一次此类地震，远多于历史记录；要么即将发生
更大的地震，要么亏损的一部分以其他方式释放。
\end{solution}

\begin{solution}{exo:b2:himalaya-andes:10}
(a) 取阿尔蒂普拉诺地壳底部的补偿深度，$61.5 - 3.7 =
\qty{57.8}{km}$，即海平面以下。普纳地壳达到 $50.5 - 4.4 = \qty{46.1}{km}$，因此
普纳柱在该深度以上含 \qty{11.7}{km} 厚的地幔：$2800\times61.5 =
\rho_P\times50.5 + 3300\times11.7$. (b) $\rho_P = (172\,200 - 38\,610)/50.5 =
\qty{2645}{kg.m^{-3}}$。(c) 两者之差 \qty{155}{kg.m^{-3}} 位于
所引 150 至 \qty{200}{kg.m^{-3}} 的下端：若仅靠地壳，则需要一个异常轻的普纳地壳。
其下较轻的岩柱更易做到：地幔岩石圈因拆沉而被移除，
代之以热而轻的软流圈。
\end{solution}

\begin{solution}{exo:b2:himalaya-andes:11}
(a) $g = 105/\tan\ang{30} = \qty{181.9}{km}$. (b) $g = 90/\tan\ang{30} + 15/\tan\ang{4} =
155.9 + 214.5 = \qty{370.4}{km}$：火山弧向陆内移动 \qty{188.5}{km}。(c) 板片
在这 \qty{300}{km} 的距离内从未到达板块之下 \qty{105}{km} 深处，
其上也没有热的地幔楔，因此不生成\omterm{def:b2:magma-emplacement:magma}{岩浆}，弧便消亡，直到板片在更东处再次下沉。
\end{solution}

\begin{solution}{exo:b2:himalaya-andes:12}
(a) \qty{21}{mm} 每年即 \qty{21}{km} 每百万年：$21\times37.5 =
\qty{787.5}{km^3}$，每千米山脉、每百万年。(b) $300/787.5 =
\qty{38}{\percent}$。(c) 体现在山脉及藏南的增厚中，
随印度一同被带到西藏之下的地壳中，以及 MHT 向北输送的物质中。
\end{solution}

\begin{solution}{pb:b2:himalaya-andes:1}
\textbf{1.} 自 \ang{28.0} 至 北纬 \ang{37.0} 的七个样品：$34\,550/7 = \qty{4936}{m}$，$h_T
= \qty{4.9}{km}$.
\textbf{2.} 自 \ang{-69} 至 \ang{-65} 的五个样品：$18\,650/5 = \qty{3730}{m}$，$h_A =
\qty{3.7}{km}$，即所引阿尔蒂普拉诺的数值；位于更南方的普纳则更高。
\textbf{3.} 西藏：至少 $9\times111.19 = \qty{1001}{km}$；阿尔蒂普拉诺：至少 $4\times
104.5 = \qty{418}{km}$.
\textbf{4.} $(5000 - 70)/(1.5\times111\,190) = 0.0296$: \qty{3.0}{\percent}.
\textbf{5.} 安第斯剖面起于海底与海沟中，低于
海平面数千米，位于高原以西；西藏剖面起于近海平面的大陆平原，
止于高耸的塔里木盆地。
\textbf{6.} $r = 4.5\times4.9 = \qty{22.05}{km}$; $T_T = 35 + 4.9 + 22.05 = \qty{62.0}{km}$.
\textbf{7.} $r = 4.5\times3.7 = \qty{16.65}{km}$; $T_A = 35 + 3.7 + 16.65 = \qty{55.4}{km}$.
\textbf{8.} 西藏：\qty{62.0}{km}，对比 70 至 \qty{75}{km}，
比中值 \qty{72.5}{km} 少 \qty{10.5}{km}（\qty{14}{\percent}）；阿尔蒂普拉诺：\qty{55.4}{km}，对比 59 至
\qty{64}{km}，比 \qty{61.5}{km} 少 \qty{6.1}{km}（\qty{10}{\percent}）。西藏偏离更远。
\textbf{9.} $72.5 - 35 = 4.9\times3300/(3300 - \rho_c)$: $3300 - \rho_c = 16\,170/37.5 =
431.2$, $\rho_c = \qty{2869}{kg.m^{-3}}$.
\textbf{10.} $26.5 = 3.7\times3300/(3300 - \rho_c)$: $3300 - \rho_c = 460.8$, $\rho_c =
\qty{2839}{kg.m^{-3}}$.
\textbf{11.} 部分荷载由板块的弯曲承担，而非局部根部；
正常地壳的厚度未必是 \qty{35}{km}；其下的地幔可能比假设的更轻或更重，
地震意义上的莫霍面未必是密度意义上地壳的底部。
\textbf{12.} 西藏地壳底部位于 $72.5 - 4.9 = \qty{67.6}{km}$，印度则位于
\qty{37.5}{km}: $t = \qty{30.1}{km}$. $\Delta g = 2\pi\times6.674\times10^{-11}\times370\times
30\,100 = \qty{4.67e-3}{m.s^{-2}}$，即相对印度 \qty{-467}{mGal}，接近
观测值 \qty{460}{mGal}。
\textbf{13.} $W_0 = 440\times60/35 = \qty{754.3}{km}$; $S = \qty{314.3}{km}$.
\textbf{14.} $(314.3 - 250)/314.3 = \qty{20}{\percent}$，即在中
安第斯所发现的缩短不足。
\textbf{15.} $0.55\times314.3 = \qty{173}{km}$，少于\omterm{def:b2:alpine-orogen:balanced}{平衡剖面}的 \qty{250}{km}，
则\omterm{def:b2:alpine-orogen:balanced}{平衡剖面}所给的数值便绰绰有余。
\textbf{16.} \omterm{def:b2:alpine-orogen:decollement}{滑脱面}之上盖层所缩短的长度，必须
与其下更靠西的基底的缩短量相匹配，因为两者在整个\omterm{def:b2:alpine-orogen:orogen}{造山带}上损失的长度相同；基底的
缩短使东科迪勒拉及高原之下的\omterm{def:b2:himalaya-andes:thickening}{地壳增厚}。
\textbf{17.} $314.3/30 = \qty{10.5}{mm}$ 每年，$10.5/72 = \qty{15}{\percent}$，占
汇聚速率的比例。
\textbf{18.} 它通过俯冲被消耗：纳斯卡板块俯冲至大陆之下。
\textbf{19.} $r = \qty{22.05}{km}$; $\Delta F = 2700\times9.8\times4900\times(35\,000 + 2450
+ 11\,025) = \qty{6.28e12}{N.m^{-1}}$.
\textbf{20.} $2700\times9.8\times3700\times45\,175 = \qty{4.42e12}{N.m^{-1}}$；比值为
0.70.
\textbf{21.} $6.28\times10^{12}/10^{5} = \qty{6.3e7}{Pa}$，约 \qty{63}{MPa}。
\textbf{22.} $9\times10^{-3}/1.06\times10^{6} = \num{8.5e-9}$ 每年。
\textbf{23.} 印度的推力沿南北向作用，并在高原内被
高原自身的重量所平衡：沿南北向高原已无法再增厚，因而竖直应力，
即重量，为最大；最小应力沿东西向，指向亚洲的自由边缘，
高原便沿此方向伸展，形成南北向的裂谷。
\textbf{24.} \textbf{\qty[text-series-to-math=true]{62}{km}}，约为
实测 \qty{72.5}{km} 的 \qty{86}{\percent}。
\end{solution}
